% =============================================================================
% NeurIPS Paper Checklist (Updated — comprehensive and honest)
% =============================================================================
% Self-contained: % =============================================================================
% NeurIPS Paper Checklist (Updated — comprehensive and honest)
% =============================================================================
% Self-contained: % =============================================================================
% NeurIPS Paper Checklist (Updated — comprehensive and honest)
% =============================================================================
% Self-contained: % =============================================================================
% NeurIPS Paper Checklist (Updated — comprehensive and honest)
% =============================================================================
% Self-contained: \input{neurips_checklist_update} into main.tex
% Replaces the section from line 1086 onward (before \end{document}).
% =============================================================================

\section*{NeurIPS Paper Checklist}

\begin{enumerate}

% --------------------------------------------------------------------------
\item {\bf Claims}
    \begin{enumerate}
        \item Do the main claims made in the abstract and introduction
              accurately reflect the paper's contributions and scope?
        \answerYes{The abstract claims that \ours{} (i) provides a unified
        multi-platform benchmark for RL post-training, (ii) reveals
        significant platform-dependent variance, and (iii) demonstrates that
        Modal/open-source pipelines are reproducible while Tinker API
        experiments are \emph{not} fully reproducible due to the API's
        closed, serverless nature.  Section~\ref{sec:results} presents
        empirical evidence for each claim.  The scope is explicitly limited
        to LoRA fine-tuning of models up to 32B parameters on three task
        domains (math, code, tool-use).  We note in
        Section~\ref{sec:limitations} that 11 of 14 Tinker experiments
        encountered JWT authentication failures and 4 of 6 Modal experiments
        encountered timeout or gradient-norm bugs; successful results are
        therefore drawn from a subset of planned runs.}
    \end{enumerate}

% --------------------------------------------------------------------------
\item {\bf Limitations}
    \begin{enumerate}
        \item Does the paper discuss the limitations of the work performed
              by the authors?
        \answerYes{Section~\ref{sec:limitations} discusses the following
        limitations explicitly:
        \begin{itemize}
          \item \textbf{Platform opacity.} Tinker API is a closed,
                serverless platform; GPU type, CUDA version, and container
                runtime are unspecified.  Results from Tinker experiments
                cannot be independently reproduced by readers without Tinker
                API access.
          \item \textbf{High experiment failure rate.} 11 of 14 Tinker
                experiments (79\%) failed due to JWT token expiry; 4 of 6
                Modal experiments (67\%) failed due to job timeouts or
                gradient norm explosions.  Statistical conclusions are
                therefore drawn from a smaller sample than planned.
          \item \textbf{Statistical power.} 5-seed TRL results and
                bootstrap confidence intervals via \texttt{rliable} are
                currently being computed; a subset of presented figures use
                preliminary estimates.
          \item \textbf{LoRA-only evaluation.} Full fine-tuning and
                quantisation-aware training are not evaluated.
          \item \textbf{Model-scale and task-coverage.} Experiments cover
                models up to 32B parameters and three task families;
                generalisation to frontier-scale (100B+) models and other
                domains is not established.
        \end{itemize}}
    \end{enumerate}

% --------------------------------------------------------------------------
\item {\bf Theory Assumptions and Proofs}
    \begin{enumerate}
        \item For each theoretical result, does the paper provide the full
              set of assumptions and a complete (or correct) proof?
        \answerNA{This is an empirical benchmarking paper; no theoretical
        results are claimed.}
    \end{enumerate}

% --------------------------------------------------------------------------
\item {\bf Experimental Result Reproducibility}
    \begin{enumerate}
        \item Does the paper fully disclose all the information needed to
              reproduce the main experimental results of the paper to the
              extent that it affects the main claims and/or conclusions of
              the paper?
        \answerPartial{Reproducibility is \emph{partial by design} due to
        platform heterogeneity.
        \begin{itemize}
          \item \textbf{Modal experiments (reproducible).}  Full code,
                Docker environments (pinned base images), seed management,
                and step-by-step commands are provided in
                \texttt{REPRODUCE.md} in the public repository
                (\url{https://github.com/arvindcr4/tinker-rl-lab}).
                Modal jobs run on explicitly provisioned NVIDIA H100 SXM5
                GPUs (80 GB) via containerised workers; all hyperparameters
                are version-controlled.  These experiments use 5 seeds with
                mean $\pm$ SE and 95\% bootstrap CIs.
          \item \textbf{Tinker API experiments (not fully reproducible).}
                Tinker is a closed-source API with serverless GPU dispatch.
                GPU type, driver version, and scheduling policy are not
                exposed.  We provide our exact API call scripts and
                configuration files, but independent replication requires a
                Tinker API account and is subject to the same JWT expiry
                issues we encountered.  We flag all Tinker-derived results
                in figures with a ``\dag{}'' symbol.
        \end{itemize}}
    \end{enumerate}

% --------------------------------------------------------------------------
\item {\bf Open access to data and code}
    \begin{enumerate}
        \item Does the paper provide open access to the data and code, with
              sufficient instructions to faithfully reproduce the main
              experimental results, as described above?
        \answerYes{All code is publicly available under MIT licence at
        \url{https://github.com/arvindcr4/tinker-rl-lab}.  Model
        checkpoints trained on Modal are uploaded to HuggingFace Hub at
        \url{https://huggingface.co/arvindcr4/tinker-rl-bench-*}.  All
        experiment runs are logged to Weights \& Biases at
        \url{https://wandb.ai/tinker-rl-lab/tinker-rl-lab-world-class};
        logs are public and include per-step reward, loss, and KL metrics.
        Training datasets (GSM8K, HumanEval, synthetic tool-use) are fully
        public and cited with licence information.  Tinker API credentials
        are \emph{not} included for security reasons; the
        \texttt{README.md} explains how to obtain a Tinker API key.}
    \end{enumerate}

% --------------------------------------------------------------------------
\item {\bf Experimental Setting/Details}
    \begin{enumerate}
        \item Does the paper specify all the training and test details
              (e.g., data splits, hyperparameters, how they were chosen,
              type of optimizer) necessary to understand the results?
        \answerYes{Section~\ref{sec:setup} and Table~\ref{tab:hyperparams}
        provide full hyperparameter settings for Modal experiments.
        Appendix~\ref{app:hyperparams} gives sweep ranges.  Tinker API
        hyperparameters are provided in the API call scripts in the
        repository, but hardware details are outside our control and are
        noted as unknown.  For Modal experiments, the GPU type (H100 SXM5,
        80 GB), batch size, learning rate schedule, LoRA rank/alpha, and
        training duration are fully specified.  Older Modal experiments used
        NVIDIA L4 GPUs (24 GB) as noted in Appendix~\ref{app:compute}.}
    \end{enumerate}

% --------------------------------------------------------------------------
\item {\bf Experiment Statistical Significance}
    \begin{enumerate}
        \item Does the paper report error bars suitably and correctly
              defined or other appropriate information about the statistical
              significance of the experiments?
        \answerPartial{For Modal/TRL experiments with 5 seeds, we report
        mean $\pm$ SE with 95\% bootstrap confidence intervals via
        \texttt{rliable}~\cite{agarwal2021deep}.  Statistical analysis is
        in progress; preliminary CIs are shown where seeds have completed.
        Tinker API experiments represent \emph{single runs} (due to high
        failure rates) and are explicitly labelled as such; no statistical
        significance is claimed for Tinker-only comparisons.  We follow the
        recommendations of~\citet{colas2019hitchhiker}
        and~\citet{jordan2024benchmarking} regarding the limits of
        single-seed evaluation.}
    \end{enumerate}

% --------------------------------------------------------------------------
\item {\bf Experiments Compute Resources}
    \begin{enumerate}
        \item For each experiment, does the paper provide sufficient
              information on the computer resources (type of compute
              workers, memory, time of execution) needed to reproduce the
              experiments?
        \answerPartial{Appendix~\ref{app:compute} lists GPU types and
        estimated GPU-hours for each experiment group.  The complete
        per-experiment breakdown, including wall-clock runtimes and
        estimated cost, is in \texttt{COMPUTE.md} in the repository.
        \begin{itemize}
          \item \textbf{Modal experiments.}  Hardware is fully specified:
                H100 SXM5 (80 GB) for recent runs, L4 (24 GB) for earlier
                experiments.  Wall-clock runtimes and GPU-hour estimates are
                provided.
          \item \textbf{Tinker API experiments.}  The Tinker platform uses
                serverless GPU dispatch; exact GPU type is not disclosed by
                the API.  Estimated GPU-hours are inferred from billing
                records and are labelled as approximate.
        \end{itemize}}
    \end{enumerate}

% --------------------------------------------------------------------------
\item {\bf Code Of Ethics}
    \begin{enumerate}
        \item Does the research conducted in the paper conform, in every
              respect, with the NeurIPS Code of Ethics?
        \answerYes{This work involves no human subjects, no private or
        sensitive data, and no models trained on proprietary corpora.  All
        training datasets are publicly released research datasets (GSM8K,
        HumanEval, synthetic tool-use data).  Broader societal impacts,
        including risks of reward hacking and alignment failure in RLHF
        systems, are discussed in the Broader Impact statement
        (Section~\ref{sec:impact} and \texttt{ethics\_statement.tex}).}
    \end{enumerate}

% --------------------------------------------------------------------------
\item {\bf Broader Impacts}
    \begin{enumerate}
        \item Does the paper discuss both potential positive societal
              impacts and negative societal impacts of the work performed?
        \answerYes{Section~\ref{sec:impact} (and the accompanying
        \texttt{ethics\_statement.tex}) discusses:
        \begin{itemize}
          \item \textbf{Positive.}  Lowering the barrier to RL post-training
                research; enabling reproducibility audits; surfacing
                platform-dependent variance that is often hidden in
                single-platform papers.
          \item \textbf{Negative / risks.}  RLHF reward hacking; alignment
                concerns when optimising proxy rewards; potential misuse of
                fine-tuned models; compute-access disparities that may limit
                replication to well-resourced groups.
          \item \textbf{Carbon footprint.}  Estimated 0.4--1.2 kg CO$_2$-eq
                for Modal H100 experiments; Tinker API footprint is
                unestimable due to undisclosed hardware.
        \end{itemize}}
    \end{enumerate}

% --------------------------------------------------------------------------
\item {\bf Safeguards}
    \begin{enumerate}
        \item Does the paper describe safeguards that have been put in place
              for responsible release of data or models that have been
              identified as requiring safeguards?
        \answerNA{The benchmark evaluates existing public models on
        mathematical reasoning, code generation, and synthetic tool-use
        tasks.  No new capabilities with clear misuse potential are
        introduced; the fine-tuned models do not exhibit harmful behaviours
        beyond their base models.  Released checkpoints are hosted on
        HuggingFace Hub under MIT licence with standard usage disclaimers.}
    \end{enumerate}

% --------------------------------------------------------------------------
\item {\bf Licenses for existing assets}
    \begin{enumerate}
        \item Are the creators or original owners of assets (e.g., code,
              data, models) used in the paper properly credited, and are
              the licence and terms of use explicitly mentioned and
              properly respected?
        \answerYes{All third-party assets are properly credited:
        \begin{itemize}
          \item \textbf{GSM8K}~\cite{cobbe2021training}: MIT licence.
          \item \textbf{HumanEval}: MIT licence (OpenAI).
          \item \textbf{Synthetic tool-use data}: self-generated from public
                API documentation; no licence restrictions.
          \item \textbf{Qwen model family}: Apache 2.0.
          \item \textbf{Llama model family}: Meta Community Licence.
          \item \textbf{TRL, PEFT, Transformers}: Apache 2.0
                (HuggingFace).
          \item \textbf{Modal}: commercial platform; terms of service
                respected.
          \item \textbf{Tinker API}: proprietary; used under standard API
                terms of service.
        \end{itemize}}
    \end{enumerate}

% --------------------------------------------------------------------------
\item {\bf New Assets}
    \begin{enumerate}
        \item Are new assets introduced in the paper well documented and is
              the documentation provided alongside the assets?
        \answerYes{New assets include: (i) the \ours{} benchmark suite
        (code, harness, evaluation scripts) released under MIT licence at
        \url{https://github.com/arvindcr4/tinker-rl-lab} with
        \texttt{README.md}, \texttt{REPRODUCE.md}, \texttt{ARTIFACT.md},
        \texttt{COMPUTE.md}, \texttt{BASELINES.md}, and
        \texttt{BENCHMARKS\_COMPARISON.md}; (ii) fine-tuned model
        checkpoints from successful Modal experiments hosted on HuggingFace
        Hub (\url{https://huggingface.co/arvindcr4/tinker-rl-bench-*});
        (iii) all experiment logs on Weights \& Biases (public project).
        Checkpoints from failed Tinker runs are \emph{not} uploaded because
        those jobs did not produce valid final states.}
    \end{enumerate}

% --------------------------------------------------------------------------
\item {\bf Crowdsourcing and Research with Human Subjects}
    \begin{enumerate}
        \item For crowdsourcing experiments and research with human
              subjects, does the paper include the full text of instructions
              given to participants and screenshots, if applicable, as well
              as details about compensation?
        \answerNA{No crowdsourcing or human subjects research was
        conducted.}
    \end{enumerate}

% --------------------------------------------------------------------------
\item {\bf Institutional Review Board (IRB) Approvals or Equivalent for
          Research with Human Subjects}
    \begin{enumerate}
        \item Does the paper describe potential risks incurred by study
              participants, whether compensation was provided to study
              participants, and whether informed consent was obtained?
        \answerNA{No human subjects research was conducted; IRB approval
        is not applicable.}
    \end{enumerate}

% --------------------------------------------------------------------------
\item {\bf Declaration of LLM usage}
    \begin{enumerate}
        \item Does the paper describe the usage of LLMs if it is an
              important, original, or non-standard component of the core
              methods?
        \answerYes{LLMs (Qwen and Llama model families) are the primary
        \emph{subjects} of evaluation; their use as RL fine-tuning targets
        is described in Sections~\ref{sec:benchmark} and~\ref{sec:setup}.
        LLMs were \textbf{not} used to draft or revise this paper.
        Code-generation assistance from GitHub Copilot was used for
        boilerplate experiment scripts; no AI-generated content appears in
        the paper itself.}
    \end{enumerate}

\end{enumerate}
 into main.tex
% Replaces the section from line 1086 onward (before \end{document}).
% =============================================================================

\section*{NeurIPS Paper Checklist}

\begin{enumerate}

% --------------------------------------------------------------------------
\item {\bf Claims}
    \begin{enumerate}
        \item Do the main claims made in the abstract and introduction
              accurately reflect the paper's contributions and scope?
        \answerYes{The abstract claims that \ours{} (i) provides a unified
        multi-platform benchmark for RL post-training, (ii) reveals
        significant platform-dependent variance, and (iii) demonstrates that
        Modal/open-source pipelines are reproducible while Tinker API
        experiments are \emph{not} fully reproducible due to the API's
        closed, serverless nature.  Section~\ref{sec:results} presents
        empirical evidence for each claim.  The scope is explicitly limited
        to LoRA fine-tuning of models up to 32B parameters on three task
        domains (math, code, tool-use).  We note in
        Section~\ref{sec:limitations} that 11 of 14 Tinker experiments
        encountered JWT authentication failures and 4 of 6 Modal experiments
        encountered timeout or gradient-norm bugs; successful results are
        therefore drawn from a subset of planned runs.}
    \end{enumerate}

% --------------------------------------------------------------------------
\item {\bf Limitations}
    \begin{enumerate}
        \item Does the paper discuss the limitations of the work performed
              by the authors?
        \answerYes{Section~\ref{sec:limitations} discusses the following
        limitations explicitly:
        \begin{itemize}
          \item \textbf{Platform opacity.} Tinker API is a closed,
                serverless platform; GPU type, CUDA version, and container
                runtime are unspecified.  Results from Tinker experiments
                cannot be independently reproduced by readers without Tinker
                API access.
          \item \textbf{High experiment failure rate.} 11 of 14 Tinker
                experiments (79\%) failed due to JWT token expiry; 4 of 6
                Modal experiments (67\%) failed due to job timeouts or
                gradient norm explosions.  Statistical conclusions are
                therefore drawn from a smaller sample than planned.
          \item \textbf{Statistical power.} 5-seed TRL results and
                bootstrap confidence intervals via \texttt{rliable} are
                currently being computed; a subset of presented figures use
                preliminary estimates.
          \item \textbf{LoRA-only evaluation.} Full fine-tuning and
                quantisation-aware training are not evaluated.
          \item \textbf{Model-scale and task-coverage.} Experiments cover
                models up to 32B parameters and three task families;
                generalisation to frontier-scale (100B+) models and other
                domains is not established.
        \end{itemize}}
    \end{enumerate}

% --------------------------------------------------------------------------
\item {\bf Theory Assumptions and Proofs}
    \begin{enumerate}
        \item For each theoretical result, does the paper provide the full
              set of assumptions and a complete (or correct) proof?
        \answerNA{This is an empirical benchmarking paper; no theoretical
        results are claimed.}
    \end{enumerate}

% --------------------------------------------------------------------------
\item {\bf Experimental Result Reproducibility}
    \begin{enumerate}
        \item Does the paper fully disclose all the information needed to
              reproduce the main experimental results of the paper to the
              extent that it affects the main claims and/or conclusions of
              the paper?
        \answerPartial{Reproducibility is \emph{partial by design} due to
        platform heterogeneity.
        \begin{itemize}
          \item \textbf{Modal experiments (reproducible).}  Full code,
                Docker environments (pinned base images), seed management,
                and step-by-step commands are provided in
                \texttt{REPRODUCE.md} in the public repository
                (\url{https://github.com/arvindcr4/tinker-rl-lab}).
                Modal jobs run on explicitly provisioned NVIDIA H100 SXM5
                GPUs (80 GB) via containerised workers; all hyperparameters
                are version-controlled.  These experiments use 5 seeds with
                mean $\pm$ SE and 95\% bootstrap CIs.
          \item \textbf{Tinker API experiments (not fully reproducible).}
                Tinker is a closed-source API with serverless GPU dispatch.
                GPU type, driver version, and scheduling policy are not
                exposed.  We provide our exact API call scripts and
                configuration files, but independent replication requires a
                Tinker API account and is subject to the same JWT expiry
                issues we encountered.  We flag all Tinker-derived results
                in figures with a ``\dag{}'' symbol.
        \end{itemize}}
    \end{enumerate}

% --------------------------------------------------------------------------
\item {\bf Open access to data and code}
    \begin{enumerate}
        \item Does the paper provide open access to the data and code, with
              sufficient instructions to faithfully reproduce the main
              experimental results, as described above?
        \answerYes{All code is publicly available under MIT licence at
        \url{https://github.com/arvindcr4/tinker-rl-lab}.  Model
        checkpoints trained on Modal are uploaded to HuggingFace Hub at
        \url{https://huggingface.co/arvindcr4/tinker-rl-bench-*}.  All
        experiment runs are logged to Weights \& Biases at
        \url{https://wandb.ai/tinker-rl-lab/tinker-rl-lab-world-class};
        logs are public and include per-step reward, loss, and KL metrics.
        Training datasets (GSM8K, HumanEval, synthetic tool-use) are fully
        public and cited with licence information.  Tinker API credentials
        are \emph{not} included for security reasons; the
        \texttt{README.md} explains how to obtain a Tinker API key.}
    \end{enumerate}

% --------------------------------------------------------------------------
\item {\bf Experimental Setting/Details}
    \begin{enumerate}
        \item Does the paper specify all the training and test details
              (e.g., data splits, hyperparameters, how they were chosen,
              type of optimizer) necessary to understand the results?
        \answerYes{Section~\ref{sec:setup} and Table~\ref{tab:hyperparams}
        provide full hyperparameter settings for Modal experiments.
        Appendix~\ref{app:hyperparams} gives sweep ranges.  Tinker API
        hyperparameters are provided in the API call scripts in the
        repository, but hardware details are outside our control and are
        noted as unknown.  For Modal experiments, the GPU type (H100 SXM5,
        80 GB), batch size, learning rate schedule, LoRA rank/alpha, and
        training duration are fully specified.  Older Modal experiments used
        NVIDIA L4 GPUs (24 GB) as noted in Appendix~\ref{app:compute}.}
    \end{enumerate}

% --------------------------------------------------------------------------
\item {\bf Experiment Statistical Significance}
    \begin{enumerate}
        \item Does the paper report error bars suitably and correctly
              defined or other appropriate information about the statistical
              significance of the experiments?
        \answerPartial{For Modal/TRL experiments with 5 seeds, we report
        mean $\pm$ SE with 95\% bootstrap confidence intervals via
        \texttt{rliable}~\cite{agarwal2021deep}.  Statistical analysis is
        in progress; preliminary CIs are shown where seeds have completed.
        Tinker API experiments represent \emph{single runs} (due to high
        failure rates) and are explicitly labelled as such; no statistical
        significance is claimed for Tinker-only comparisons.  We follow the
        recommendations of~\citet{colas2019hitchhiker}
        and~\citet{jordan2024benchmarking} regarding the limits of
        single-seed evaluation.}
    \end{enumerate}

% --------------------------------------------------------------------------
\item {\bf Experiments Compute Resources}
    \begin{enumerate}
        \item For each experiment, does the paper provide sufficient
              information on the computer resources (type of compute
              workers, memory, time of execution) needed to reproduce the
              experiments?
        \answerPartial{Appendix~\ref{app:compute} lists GPU types and
        estimated GPU-hours for each experiment group.  The complete
        per-experiment breakdown, including wall-clock runtimes and
        estimated cost, is in \texttt{COMPUTE.md} in the repository.
        \begin{itemize}
          \item \textbf{Modal experiments.}  Hardware is fully specified:
                H100 SXM5 (80 GB) for recent runs, L4 (24 GB) for earlier
                experiments.  Wall-clock runtimes and GPU-hour estimates are
                provided.
          \item \textbf{Tinker API experiments.}  The Tinker platform uses
                serverless GPU dispatch; exact GPU type is not disclosed by
                the API.  Estimated GPU-hours are inferred from billing
                records and are labelled as approximate.
        \end{itemize}}
    \end{enumerate}

% --------------------------------------------------------------------------
\item {\bf Code Of Ethics}
    \begin{enumerate}
        \item Does the research conducted in the paper conform, in every
              respect, with the NeurIPS Code of Ethics?
        \answerYes{This work involves no human subjects, no private or
        sensitive data, and no models trained on proprietary corpora.  All
        training datasets are publicly released research datasets (GSM8K,
        HumanEval, synthetic tool-use data).  Broader societal impacts,
        including risks of reward hacking and alignment failure in RLHF
        systems, are discussed in the Broader Impact statement
        (Section~\ref{sec:impact} and \texttt{ethics\_statement.tex}).}
    \end{enumerate}

% --------------------------------------------------------------------------
\item {\bf Broader Impacts}
    \begin{enumerate}
        \item Does the paper discuss both potential positive societal
              impacts and negative societal impacts of the work performed?
        \answerYes{Section~\ref{sec:impact} (and the accompanying
        \texttt{ethics\_statement.tex}) discusses:
        \begin{itemize}
          \item \textbf{Positive.}  Lowering the barrier to RL post-training
                research; enabling reproducibility audits; surfacing
                platform-dependent variance that is often hidden in
                single-platform papers.
          \item \textbf{Negative / risks.}  RLHF reward hacking; alignment
                concerns when optimising proxy rewards; potential misuse of
                fine-tuned models; compute-access disparities that may limit
                replication to well-resourced groups.
          \item \textbf{Carbon footprint.}  Estimated 0.4--1.2 kg CO$_2$-eq
                for Modal H100 experiments; Tinker API footprint is
                unestimable due to undisclosed hardware.
        \end{itemize}}
    \end{enumerate}

% --------------------------------------------------------------------------
\item {\bf Safeguards}
    \begin{enumerate}
        \item Does the paper describe safeguards that have been put in place
              for responsible release of data or models that have been
              identified as requiring safeguards?
        \answerNA{The benchmark evaluates existing public models on
        mathematical reasoning, code generation, and synthetic tool-use
        tasks.  No new capabilities with clear misuse potential are
        introduced; the fine-tuned models do not exhibit harmful behaviours
        beyond their base models.  Released checkpoints are hosted on
        HuggingFace Hub under MIT licence with standard usage disclaimers.}
    \end{enumerate}

% --------------------------------------------------------------------------
\item {\bf Licenses for existing assets}
    \begin{enumerate}
        \item Are the creators or original owners of assets (e.g., code,
              data, models) used in the paper properly credited, and are
              the licence and terms of use explicitly mentioned and
              properly respected?
        \answerYes{All third-party assets are properly credited:
        \begin{itemize}
          \item \textbf{GSM8K}~\cite{cobbe2021training}: MIT licence.
          \item \textbf{HumanEval}: MIT licence (OpenAI).
          \item \textbf{Synthetic tool-use data}: self-generated from public
                API documentation; no licence restrictions.
          \item \textbf{Qwen model family}: Apache 2.0.
          \item \textbf{Llama model family}: Meta Community Licence.
          \item \textbf{TRL, PEFT, Transformers}: Apache 2.0
                (HuggingFace).
          \item \textbf{Modal}: commercial platform; terms of service
                respected.
          \item \textbf{Tinker API}: proprietary; used under standard API
                terms of service.
        \end{itemize}}
    \end{enumerate}

% --------------------------------------------------------------------------
\item {\bf New Assets}
    \begin{enumerate}
        \item Are new assets introduced in the paper well documented and is
              the documentation provided alongside the assets?
        \answerYes{New assets include: (i) the \ours{} benchmark suite
        (code, harness, evaluation scripts) released under MIT licence at
        \url{https://github.com/arvindcr4/tinker-rl-lab} with
        \texttt{README.md}, \texttt{REPRODUCE.md}, \texttt{ARTIFACT.md},
        \texttt{COMPUTE.md}, \texttt{BASELINES.md}, and
        \texttt{BENCHMARKS\_COMPARISON.md}; (ii) fine-tuned model
        checkpoints from successful Modal experiments hosted on HuggingFace
        Hub (\url{https://huggingface.co/arvindcr4/tinker-rl-bench-*});
        (iii) all experiment logs on Weights \& Biases (public project).
        Checkpoints from failed Tinker runs are \emph{not} uploaded because
        those jobs did not produce valid final states.}
    \end{enumerate}

% --------------------------------------------------------------------------
\item {\bf Crowdsourcing and Research with Human Subjects}
    \begin{enumerate}
        \item For crowdsourcing experiments and research with human
              subjects, does the paper include the full text of instructions
              given to participants and screenshots, if applicable, as well
              as details about compensation?
        \answerNA{No crowdsourcing or human subjects research was
        conducted.}
    \end{enumerate}

% --------------------------------------------------------------------------
\item {\bf Institutional Review Board (IRB) Approvals or Equivalent for
          Research with Human Subjects}
    \begin{enumerate}
        \item Does the paper describe potential risks incurred by study
              participants, whether compensation was provided to study
              participants, and whether informed consent was obtained?
        \answerNA{No human subjects research was conducted; IRB approval
        is not applicable.}
    \end{enumerate}

% --------------------------------------------------------------------------
\item {\bf Declaration of LLM usage}
    \begin{enumerate}
        \item Does the paper describe the usage of LLMs if it is an
              important, original, or non-standard component of the core
              methods?
        \answerYes{LLMs (Qwen and Llama model families) are the primary
        \emph{subjects} of evaluation; their use as RL fine-tuning targets
        is described in Sections~\ref{sec:benchmark} and~\ref{sec:setup}.
        LLMs were \textbf{not} used to draft or revise this paper.
        Code-generation assistance from GitHub Copilot was used for
        boilerplate experiment scripts; no AI-generated content appears in
        the paper itself.}
    \end{enumerate}

\end{enumerate}
 into main.tex
% Replaces the section from line 1086 onward (before \end{document}).
% =============================================================================

\section*{NeurIPS Paper Checklist}

\begin{enumerate}

% --------------------------------------------------------------------------
\item {\bf Claims}
    \begin{enumerate}
        \item Do the main claims made in the abstract and introduction
              accurately reflect the paper's contributions and scope?
        \answerYes{The abstract claims that \ours{} (i) provides a unified
        multi-platform benchmark for RL post-training, (ii) reveals
        significant platform-dependent variance, and (iii) demonstrates that
        Modal/open-source pipelines are reproducible while Tinker API
        experiments are \emph{not} fully reproducible due to the API's
        closed, serverless nature.  Section~\ref{sec:results} presents
        empirical evidence for each claim.  The scope is explicitly limited
        to LoRA fine-tuning of models up to 32B parameters on three task
        domains (math, code, tool-use).  We note in
        Section~\ref{sec:limitations} that 11 of 14 Tinker experiments
        encountered JWT authentication failures and 4 of 6 Modal experiments
        encountered timeout or gradient-norm bugs; successful results are
        therefore drawn from a subset of planned runs.}
    \end{enumerate}

% --------------------------------------------------------------------------
\item {\bf Limitations}
    \begin{enumerate}
        \item Does the paper discuss the limitations of the work performed
              by the authors?
        \answerYes{Section~\ref{sec:limitations} discusses the following
        limitations explicitly:
        \begin{itemize}
          \item \textbf{Platform opacity.} Tinker API is a closed,
                serverless platform; GPU type, CUDA version, and container
                runtime are unspecified.  Results from Tinker experiments
                cannot be independently reproduced by readers without Tinker
                API access.
          \item \textbf{High experiment failure rate.} 11 of 14 Tinker
                experiments (79\%) failed due to JWT token expiry; 4 of 6
                Modal experiments (67\%) failed due to job timeouts or
                gradient norm explosions.  Statistical conclusions are
                therefore drawn from a smaller sample than planned.
          \item \textbf{Statistical power.} 5-seed TRL results and
                bootstrap confidence intervals via \texttt{rliable} are
                currently being computed; a subset of presented figures use
                preliminary estimates.
          \item \textbf{LoRA-only evaluation.} Full fine-tuning and
                quantisation-aware training are not evaluated.
          \item \textbf{Model-scale and task-coverage.} Experiments cover
                models up to 32B parameters and three task families;
                generalisation to frontier-scale (100B+) models and other
                domains is not established.
        \end{itemize}}
    \end{enumerate}

% --------------------------------------------------------------------------
\item {\bf Theory Assumptions and Proofs}
    \begin{enumerate}
        \item For each theoretical result, does the paper provide the full
              set of assumptions and a complete (or correct) proof?
        \answerNA{This is an empirical benchmarking paper; no theoretical
        results are claimed.}
    \end{enumerate}

% --------------------------------------------------------------------------
\item {\bf Experimental Result Reproducibility}
    \begin{enumerate}
        \item Does the paper fully disclose all the information needed to
              reproduce the main experimental results of the paper to the
              extent that it affects the main claims and/or conclusions of
              the paper?
        \answerPartial{Reproducibility is \emph{partial by design} due to
        platform heterogeneity.
        \begin{itemize}
          \item \textbf{Modal experiments (reproducible).}  Full code,
                Docker environments (pinned base images), seed management,
                and step-by-step commands are provided in
                \texttt{REPRODUCE.md} in the public repository
                (\url{https://github.com/arvindcr4/tinker-rl-lab}).
                Modal jobs run on explicitly provisioned NVIDIA H100 SXM5
                GPUs (80 GB) via containerised workers; all hyperparameters
                are version-controlled.  These experiments use 5 seeds with
                mean $\pm$ SE and 95\% bootstrap CIs.
          \item \textbf{Tinker API experiments (not fully reproducible).}
                Tinker is a closed-source API with serverless GPU dispatch.
                GPU type, driver version, and scheduling policy are not
                exposed.  We provide our exact API call scripts and
                configuration files, but independent replication requires a
                Tinker API account and is subject to the same JWT expiry
                issues we encountered.  We flag all Tinker-derived results
                in figures with a ``\dag{}'' symbol.
        \end{itemize}}
    \end{enumerate}

% --------------------------------------------------------------------------
\item {\bf Open access to data and code}
    \begin{enumerate}
        \item Does the paper provide open access to the data and code, with
              sufficient instructions to faithfully reproduce the main
              experimental results, as described above?
        \answerYes{All code is publicly available under MIT licence at
        \url{https://github.com/arvindcr4/tinker-rl-lab}.  Model
        checkpoints trained on Modal are uploaded to HuggingFace Hub at
        \url{https://huggingface.co/arvindcr4/tinker-rl-bench-*}.  All
        experiment runs are logged to Weights \& Biases at
        \url{https://wandb.ai/tinker-rl-lab/tinker-rl-lab-world-class};
        logs are public and include per-step reward, loss, and KL metrics.
        Training datasets (GSM8K, HumanEval, synthetic tool-use) are fully
        public and cited with licence information.  Tinker API credentials
        are \emph{not} included for security reasons; the
        \texttt{README.md} explains how to obtain a Tinker API key.}
    \end{enumerate}

% --------------------------------------------------------------------------
\item {\bf Experimental Setting/Details}
    \begin{enumerate}
        \item Does the paper specify all the training and test details
              (e.g., data splits, hyperparameters, how they were chosen,
              type of optimizer) necessary to understand the results?
        \answerYes{Section~\ref{sec:setup} and Table~\ref{tab:hyperparams}
        provide full hyperparameter settings for Modal experiments.
        Appendix~\ref{app:hyperparams} gives sweep ranges.  Tinker API
        hyperparameters are provided in the API call scripts in the
        repository, but hardware details are outside our control and are
        noted as unknown.  For Modal experiments, the GPU type (H100 SXM5,
        80 GB), batch size, learning rate schedule, LoRA rank/alpha, and
        training duration are fully specified.  Older Modal experiments used
        NVIDIA L4 GPUs (24 GB) as noted in Appendix~\ref{app:compute}.}
    \end{enumerate}

% --------------------------------------------------------------------------
\item {\bf Experiment Statistical Significance}
    \begin{enumerate}
        \item Does the paper report error bars suitably and correctly
              defined or other appropriate information about the statistical
              significance of the experiments?
        \answerPartial{For Modal/TRL experiments with 5 seeds, we report
        mean $\pm$ SE with 95\% bootstrap confidence intervals via
        \texttt{rliable}~\cite{agarwal2021deep}.  Statistical analysis is
        in progress; preliminary CIs are shown where seeds have completed.
        Tinker API experiments represent \emph{single runs} (due to high
        failure rates) and are explicitly labelled as such; no statistical
        significance is claimed for Tinker-only comparisons.  We follow the
        recommendations of~\citet{colas2019hitchhiker}
        and~\citet{jordan2024benchmarking} regarding the limits of
        single-seed evaluation.}
    \end{enumerate}

% --------------------------------------------------------------------------
\item {\bf Experiments Compute Resources}
    \begin{enumerate}
        \item For each experiment, does the paper provide sufficient
              information on the computer resources (type of compute
              workers, memory, time of execution) needed to reproduce the
              experiments?
        \answerPartial{Appendix~\ref{app:compute} lists GPU types and
        estimated GPU-hours for each experiment group.  The complete
        per-experiment breakdown, including wall-clock runtimes and
        estimated cost, is in \texttt{COMPUTE.md} in the repository.
        \begin{itemize}
          \item \textbf{Modal experiments.}  Hardware is fully specified:
                H100 SXM5 (80 GB) for recent runs, L4 (24 GB) for earlier
                experiments.  Wall-clock runtimes and GPU-hour estimates are
                provided.
          \item \textbf{Tinker API experiments.}  The Tinker platform uses
                serverless GPU dispatch; exact GPU type is not disclosed by
                the API.  Estimated GPU-hours are inferred from billing
                records and are labelled as approximate.
        \end{itemize}}
    \end{enumerate}

% --------------------------------------------------------------------------
\item {\bf Code Of Ethics}
    \begin{enumerate}
        \item Does the research conducted in the paper conform, in every
              respect, with the NeurIPS Code of Ethics?
        \answerYes{This work involves no human subjects, no private or
        sensitive data, and no models trained on proprietary corpora.  All
        training datasets are publicly released research datasets (GSM8K,
        HumanEval, synthetic tool-use data).  Broader societal impacts,
        including risks of reward hacking and alignment failure in RLHF
        systems, are discussed in the Broader Impact statement
        (Section~\ref{sec:impact} and \texttt{ethics\_statement.tex}).}
    \end{enumerate}

% --------------------------------------------------------------------------
\item {\bf Broader Impacts}
    \begin{enumerate}
        \item Does the paper discuss both potential positive societal
              impacts and negative societal impacts of the work performed?
        \answerYes{Section~\ref{sec:impact} (and the accompanying
        \texttt{ethics\_statement.tex}) discusses:
        \begin{itemize}
          \item \textbf{Positive.}  Lowering the barrier to RL post-training
                research; enabling reproducibility audits; surfacing
                platform-dependent variance that is often hidden in
                single-platform papers.
          \item \textbf{Negative / risks.}  RLHF reward hacking; alignment
                concerns when optimising proxy rewards; potential misuse of
                fine-tuned models; compute-access disparities that may limit
                replication to well-resourced groups.
          \item \textbf{Carbon footprint.}  Estimated 0.4--1.2 kg CO$_2$-eq
                for Modal H100 experiments; Tinker API footprint is
                unestimable due to undisclosed hardware.
        \end{itemize}}
    \end{enumerate}

% --------------------------------------------------------------------------
\item {\bf Safeguards}
    \begin{enumerate}
        \item Does the paper describe safeguards that have been put in place
              for responsible release of data or models that have been
              identified as requiring safeguards?
        \answerNA{The benchmark evaluates existing public models on
        mathematical reasoning, code generation, and synthetic tool-use
        tasks.  No new capabilities with clear misuse potential are
        introduced; the fine-tuned models do not exhibit harmful behaviours
        beyond their base models.  Released checkpoints are hosted on
        HuggingFace Hub under MIT licence with standard usage disclaimers.}
    \end{enumerate}

% --------------------------------------------------------------------------
\item {\bf Licenses for existing assets}
    \begin{enumerate}
        \item Are the creators or original owners of assets (e.g., code,
              data, models) used in the paper properly credited, and are
              the licence and terms of use explicitly mentioned and
              properly respected?
        \answerYes{All third-party assets are properly credited:
        \begin{itemize}
          \item \textbf{GSM8K}~\cite{cobbe2021training}: MIT licence.
          \item \textbf{HumanEval}: MIT licence (OpenAI).
          \item \textbf{Synthetic tool-use data}: self-generated from public
                API documentation; no licence restrictions.
          \item \textbf{Qwen model family}: Apache 2.0.
          \item \textbf{Llama model family}: Meta Community Licence.
          \item \textbf{TRL, PEFT, Transformers}: Apache 2.0
                (HuggingFace).
          \item \textbf{Modal}: commercial platform; terms of service
                respected.
          \item \textbf{Tinker API}: proprietary; used under standard API
                terms of service.
        \end{itemize}}
    \end{enumerate}

% --------------------------------------------------------------------------
\item {\bf New Assets}
    \begin{enumerate}
        \item Are new assets introduced in the paper well documented and is
              the documentation provided alongside the assets?
        \answerYes{New assets include: (i) the \ours{} benchmark suite
        (code, harness, evaluation scripts) released under MIT licence at
        \url{https://github.com/arvindcr4/tinker-rl-lab} with
        \texttt{README.md}, \texttt{REPRODUCE.md}, \texttt{ARTIFACT.md},
        \texttt{COMPUTE.md}, \texttt{BASELINES.md}, and
        \texttt{BENCHMARKS\_COMPARISON.md}; (ii) fine-tuned model
        checkpoints from successful Modal experiments hosted on HuggingFace
        Hub (\url{https://huggingface.co/arvindcr4/tinker-rl-bench-*});
        (iii) all experiment logs on Weights \& Biases (public project).
        Checkpoints from failed Tinker runs are \emph{not} uploaded because
        those jobs did not produce valid final states.}
    \end{enumerate}

% --------------------------------------------------------------------------
\item {\bf Crowdsourcing and Research with Human Subjects}
    \begin{enumerate}
        \item For crowdsourcing experiments and research with human
              subjects, does the paper include the full text of instructions
              given to participants and screenshots, if applicable, as well
              as details about compensation?
        \answerNA{No crowdsourcing or human subjects research was
        conducted.}
    \end{enumerate}

% --------------------------------------------------------------------------
\item {\bf Institutional Review Board (IRB) Approvals or Equivalent for
          Research with Human Subjects}
    \begin{enumerate}
        \item Does the paper describe potential risks incurred by study
              participants, whether compensation was provided to study
              participants, and whether informed consent was obtained?
        \answerNA{No human subjects research was conducted; IRB approval
        is not applicable.}
    \end{enumerate}

% --------------------------------------------------------------------------
\item {\bf Declaration of LLM usage}
    \begin{enumerate}
        \item Does the paper describe the usage of LLMs if it is an
              important, original, or non-standard component of the core
              methods?
        \answerYes{LLMs (Qwen and Llama model families) are the primary
        \emph{subjects} of evaluation; their use as RL fine-tuning targets
        is described in Sections~\ref{sec:benchmark} and~\ref{sec:setup}.
        LLMs were \textbf{not} used to draft or revise this paper.
        Code-generation assistance from GitHub Copilot was used for
        boilerplate experiment scripts; no AI-generated content appears in
        the paper itself.}
    \end{enumerate}

\end{enumerate}
 into main.tex
% Replaces the section from line 1086 onward (before \end{document}).
% =============================================================================

\section*{NeurIPS Paper Checklist}

\begin{enumerate}

% --------------------------------------------------------------------------
\item {\bf Claims}
    \begin{enumerate}
        \item Do the main claims made in the abstract and introduction
              accurately reflect the paper's contributions and scope?
        \answerYes{The abstract claims that \ours{} (i) provides a unified
        multi-platform benchmark for RL post-training, (ii) reveals
        significant platform-dependent variance, and (iii) demonstrates that
        Modal/open-source pipelines are reproducible while Tinker API
        experiments are \emph{not} fully reproducible due to the API's
        closed, serverless nature.  Section~\ref{sec:results} presents
        empirical evidence for each claim.  The scope is explicitly limited
        to LoRA fine-tuning of models up to 32B parameters on three task
        domains (math, code, tool-use).  We note in
        Section~\ref{sec:limitations} that 11 of 14 Tinker experiments
        encountered JWT authentication failures and 4 of 6 Modal experiments
        encountered timeout or gradient-norm bugs; successful results are
        therefore drawn from a subset of planned runs.}
    \end{enumerate}

% --------------------------------------------------------------------------
\item {\bf Limitations}
    \begin{enumerate}
        \item Does the paper discuss the limitations of the work performed
              by the authors?
        \answerYes{Section~\ref{sec:limitations} discusses the following
        limitations explicitly:
        \begin{itemize}
          \item \textbf{Platform opacity.} Tinker API is a closed,
                serverless platform; GPU type, CUDA version, and container
                runtime are unspecified.  Results from Tinker experiments
                cannot be independently reproduced by readers without Tinker
                API access.
          \item \textbf{High experiment failure rate.} 11 of 14 Tinker
                experiments (79\%) failed due to JWT token expiry; 4 of 6
                Modal experiments (67\%) failed due to job timeouts or
                gradient norm explosions.  Statistical conclusions are
                therefore drawn from a smaller sample than planned.
          \item \textbf{Statistical power.} 5-seed TRL results and
                bootstrap confidence intervals via \texttt{rliable} are
                currently being computed; a subset of presented figures use
                preliminary estimates.
          \item \textbf{LoRA-only evaluation.} Full fine-tuning and
                quantisation-aware training are not evaluated.
          \item \textbf{Model-scale and task-coverage.} Experiments cover
                models up to 32B parameters and three task families;
                generalisation to frontier-scale (100B+) models and other
                domains is not established.
        \end{itemize}}
    \end{enumerate}

% --------------------------------------------------------------------------
\item {\bf Theory Assumptions and Proofs}
    \begin{enumerate}
        \item For each theoretical result, does the paper provide the full
              set of assumptions and a complete (or correct) proof?
        \answerNA{This is an empirical benchmarking paper; no theoretical
        results are claimed.}
    \end{enumerate}

% --------------------------------------------------------------------------
\item {\bf Experimental Result Reproducibility}
    \begin{enumerate}
        \item Does the paper fully disclose all the information needed to
              reproduce the main experimental results of the paper to the
              extent that it affects the main claims and/or conclusions of
              the paper?
        \answerPartial{Reproducibility is \emph{partial by design} due to
        platform heterogeneity.
        \begin{itemize}
          \item \textbf{Modal experiments (reproducible).}  Full code,
                Docker environments (pinned base images), seed management,
                and step-by-step commands are provided in
                \texttt{REPRODUCE.md} in the public repository
                (\url{https://github.com/arvindcr4/tinker-rl-lab}).
                Modal jobs run on explicitly provisioned NVIDIA H100 SXM5
                GPUs (80 GB) via containerised workers; all hyperparameters
                are version-controlled.  These experiments use 5 seeds with
                mean $\pm$ SE and 95\% bootstrap CIs.
          \item \textbf{Tinker API experiments (not fully reproducible).}
                Tinker is a closed-source API with serverless GPU dispatch.
                GPU type, driver version, and scheduling policy are not
                exposed.  We provide our exact API call scripts and
                configuration files, but independent replication requires a
                Tinker API account and is subject to the same JWT expiry
                issues we encountered.  We flag all Tinker-derived results
                in figures with a ``\dag{}'' symbol.
        \end{itemize}}
    \end{enumerate}

% --------------------------------------------------------------------------
\item {\bf Open access to data and code}
    \begin{enumerate}
        \item Does the paper provide open access to the data and code, with
              sufficient instructions to faithfully reproduce the main
              experimental results, as described above?
        \answerYes{All code is publicly available under MIT licence at
        \url{https://github.com/arvindcr4/tinker-rl-lab}.  Model
        checkpoints trained on Modal are uploaded to HuggingFace Hub at
        \url{https://huggingface.co/arvindcr4/tinker-rl-bench-*}.  All
        experiment runs are logged to Weights \& Biases at
        \url{https://wandb.ai/tinker-rl-lab/tinker-rl-lab-world-class};
        logs are public and include per-step reward, loss, and KL metrics.
        Training datasets (GSM8K, HumanEval, synthetic tool-use) are fully
        public and cited with licence information.  Tinker API credentials
        are \emph{not} included for security reasons; the
        \texttt{README.md} explains how to obtain a Tinker API key.}
    \end{enumerate}

% --------------------------------------------------------------------------
\item {\bf Experimental Setting/Details}
    \begin{enumerate}
        \item Does the paper specify all the training and test details
              (e.g., data splits, hyperparameters, how they were chosen,
              type of optimizer) necessary to understand the results?
        \answerYes{Section~\ref{sec:setup} and Table~\ref{tab:hyperparams}
        provide full hyperparameter settings for Modal experiments.
        Appendix~\ref{app:hyperparams} gives sweep ranges.  Tinker API
        hyperparameters are provided in the API call scripts in the
        repository, but hardware details are outside our control and are
        noted as unknown.  For Modal experiments, the GPU type (H100 SXM5,
        80 GB), batch size, learning rate schedule, LoRA rank/alpha, and
        training duration are fully specified.  Older Modal experiments used
        NVIDIA L4 GPUs (24 GB) as noted in Appendix~\ref{app:compute}.}
    \end{enumerate}

% --------------------------------------------------------------------------
\item {\bf Experiment Statistical Significance}
    \begin{enumerate}
        \item Does the paper report error bars suitably and correctly
              defined or other appropriate information about the statistical
              significance of the experiments?
        \answerPartial{For Modal/TRL experiments with 5 seeds, we report
        mean $\pm$ SE with 95\% bootstrap confidence intervals via
        \texttt{rliable}~\cite{agarwal2021deep}.  Statistical analysis is
        in progress; preliminary CIs are shown where seeds have completed.
        Tinker API experiments represent \emph{single runs} (due to high
        failure rates) and are explicitly labelled as such; no statistical
        significance is claimed for Tinker-only comparisons.  We follow the
        recommendations of~\citet{colas2019hitchhiker}
        and~\citet{jordan2024benchmarking} regarding the limits of
        single-seed evaluation.}
    \end{enumerate}

% --------------------------------------------------------------------------
\item {\bf Experiments Compute Resources}
    \begin{enumerate}
        \item For each experiment, does the paper provide sufficient
              information on the computer resources (type of compute
              workers, memory, time of execution) needed to reproduce the
              experiments?
        \answerPartial{Appendix~\ref{app:compute} lists GPU types and
        estimated GPU-hours for each experiment group.  The complete
        per-experiment breakdown, including wall-clock runtimes and
        estimated cost, is in \texttt{COMPUTE.md} in the repository.
        \begin{itemize}
          \item \textbf{Modal experiments.}  Hardware is fully specified:
                H100 SXM5 (80 GB) for recent runs, L4 (24 GB) for earlier
                experiments.  Wall-clock runtimes and GPU-hour estimates are
                provided.
          \item \textbf{Tinker API experiments.}  The Tinker platform uses
                serverless GPU dispatch; exact GPU type is not disclosed by
                the API.  Estimated GPU-hours are inferred from billing
                records and are labelled as approximate.
        \end{itemize}}
    \end{enumerate}

% --------------------------------------------------------------------------
\item {\bf Code Of Ethics}
    \begin{enumerate}
        \item Does the research conducted in the paper conform, in every
              respect, with the NeurIPS Code of Ethics?
        \answerYes{This work involves no human subjects, no private or
        sensitive data, and no models trained on proprietary corpora.  All
        training datasets are publicly released research datasets (GSM8K,
        HumanEval, synthetic tool-use data).  Broader societal impacts,
        including risks of reward hacking and alignment failure in RLHF
        systems, are discussed in the Broader Impact statement
        (Section~\ref{sec:impact} and \texttt{ethics\_statement.tex}).}
    \end{enumerate}

% --------------------------------------------------------------------------
\item {\bf Broader Impacts}
    \begin{enumerate}
        \item Does the paper discuss both potential positive societal
              impacts and negative societal impacts of the work performed?
        \answerYes{Section~\ref{sec:impact} (and the accompanying
        \texttt{ethics\_statement.tex}) discusses:
        \begin{itemize}
          \item \textbf{Positive.}  Lowering the barrier to RL post-training
                research; enabling reproducibility audits; surfacing
                platform-dependent variance that is often hidden in
                single-platform papers.
          \item \textbf{Negative / risks.}  RLHF reward hacking; alignment
                concerns when optimising proxy rewards; potential misuse of
                fine-tuned models; compute-access disparities that may limit
                replication to well-resourced groups.
          \item \textbf{Carbon footprint.}  Estimated 0.4--1.2 kg CO$_2$-eq
                for Modal H100 experiments; Tinker API footprint is
                unestimable due to undisclosed hardware.
        \end{itemize}}
    \end{enumerate}

% --------------------------------------------------------------------------
\item {\bf Safeguards}
    \begin{enumerate}
        \item Does the paper describe safeguards that have been put in place
              for responsible release of data or models that have been
              identified as requiring safeguards?
        \answerNA{The benchmark evaluates existing public models on
        mathematical reasoning, code generation, and synthetic tool-use
        tasks.  No new capabilities with clear misuse potential are
        introduced; the fine-tuned models do not exhibit harmful behaviours
        beyond their base models.  Released checkpoints are hosted on
        HuggingFace Hub under MIT licence with standard usage disclaimers.}
    \end{enumerate}

% --------------------------------------------------------------------------
\item {\bf Licenses for existing assets}
    \begin{enumerate}
        \item Are the creators or original owners of assets (e.g., code,
              data, models) used in the paper properly credited, and are
              the licence and terms of use explicitly mentioned and
              properly respected?
        \answerYes{All third-party assets are properly credited:
        \begin{itemize}
          \item \textbf{GSM8K}~\cite{cobbe2021training}: MIT licence.
          \item \textbf{HumanEval}: MIT licence (OpenAI).
          \item \textbf{Synthetic tool-use data}: self-generated from public
                API documentation; no licence restrictions.
          \item \textbf{Qwen model family}: Apache 2.0.
          \item \textbf{Llama model family}: Meta Community Licence.
          \item \textbf{TRL, PEFT, Transformers}: Apache 2.0
                (HuggingFace).
          \item \textbf{Modal}: commercial platform; terms of service
                respected.
          \item \textbf{Tinker API}: proprietary; used under standard API
                terms of service.
        \end{itemize}}
    \end{enumerate}

% --------------------------------------------------------------------------
\item {\bf New Assets}
    \begin{enumerate}
        \item Are new assets introduced in the paper well documented and is
              the documentation provided alongside the assets?
        \answerYes{New assets include: (i) the \ours{} benchmark suite
        (code, harness, evaluation scripts) released under MIT licence at
        \url{https://github.com/arvindcr4/tinker-rl-lab} with
        \texttt{README.md}, \texttt{REPRODUCE.md}, \texttt{ARTIFACT.md},
        \texttt{COMPUTE.md}, \texttt{BASELINES.md}, and
        \texttt{BENCHMARKS\_COMPARISON.md}; (ii) fine-tuned model
        checkpoints from successful Modal experiments hosted on HuggingFace
        Hub (\url{https://huggingface.co/arvindcr4/tinker-rl-bench-*});
        (iii) all experiment logs on Weights \& Biases (public project).
        Checkpoints from failed Tinker runs are \emph{not} uploaded because
        those jobs did not produce valid final states.}
    \end{enumerate}

% --------------------------------------------------------------------------
\item {\bf Crowdsourcing and Research with Human Subjects}
    \begin{enumerate}
        \item For crowdsourcing experiments and research with human
              subjects, does the paper include the full text of instructions
              given to participants and screenshots, if applicable, as well
              as details about compensation?
        \answerNA{No crowdsourcing or human subjects research was
        conducted.}
    \end{enumerate}

% --------------------------------------------------------------------------
\item {\bf Institutional Review Board (IRB) Approvals or Equivalent for
          Research with Human Subjects}
    \begin{enumerate}
        \item Does the paper describe potential risks incurred by study
              participants, whether compensation was provided to study
              participants, and whether informed consent was obtained?
        \answerNA{No human subjects research was conducted; IRB approval
        is not applicable.}
    \end{enumerate}

% --------------------------------------------------------------------------
\item {\bf Declaration of LLM usage}
    \begin{enumerate}
        \item Does the paper describe the usage of LLMs if it is an
              important, original, or non-standard component of the core
              methods?
        \answerYes{LLMs (Qwen and Llama model families) are the primary
        \emph{subjects} of evaluation; their use as RL fine-tuning targets
        is described in Sections~\ref{sec:benchmark} and~\ref{sec:setup}.
        LLMs were \textbf{not} used to draft or revise this paper.
        Code-generation assistance from GitHub Copilot was used for
        boilerplate experiment scripts; no AI-generated content appears in
        the paper itself.}
    \end{enumerate}

\end{enumerate}
